\section{Acumulación de errores y corrección en línea}

\subsection{Origen de los compounding errors}

El challenge central del aprendizaje por imitación consiste en que la acción predicha $\hat{a}$ modifica el siguiente estado $s'$, lo cual produce una distribución de entrada no i.i.d. Los errores pequeños se amplifican progresivamente, hasta producir una catastrophic failure.

\begin{warningbox}{Desplazamiento de distribución y covariate shift}
Que la distribución de estados del experto $p_{\text{expert}}(s)$ y la distribución de la política aprendida $p_{\pi}(s)$ no coincidan es un riesgo sistemático: si no se interviene, aun cuando la loss promedio de entrenamiento sea baja, el performance puede desplomarse en unos cuantos estados.
\end{warningbox}

\subsection{DAgger y Human-Gated DAgger}

La idea de solución consiste en dejar que la política se ejecute por sí misma, recolectar los nuevos estados y usar al experto para proporcionar acciones de corrección en esos estados:
\begin{enumerate}
    \item Ejecutar un rollout de la política actual $\pi_\theta$ y registrar los estados visitados.
    \item Solicitar la acción del experto en cada estado y agregarla al dataset.
    \item Reentrenar la política o actualizar los parámetros mediante importance weighting.
\end{enumerate}

\begin{table}[H]
\centering
\begin{tabular}{p{0.28\textwidth} p{0.28\textwidth} p{0.28\textwidth}}
\toprule
Etapa & Entregable clave & Elemento de revisión \\
\midrule
Observe & Rollout log & state distribution shift \\
Assess & Expert correction & human agreement rate \\
Act & Updated policy checkpoint & performance vs baseline \\
\bottomrule
\end{tabular}
\caption{Entregables y revisiones en el proceso de DAgger}
\end{table}

\begin{knowledgebox}{Offline vs. Online del aprendizaje por imitación}
\begin{itemize}
    \item \textbf{Offline (clonación de comportamiento)}: solo usa un static dataset; los datos son seguros, pero la generalización es deficiente.
    \item \textbf{Online (DAgger/HG-DAgger)}: implica human-in-the-loop; es eficiente, pero requiere permisos de experto.
\end{itemize}
\end{knowledgebox}

\subsection{Sinergia entre el aprendizaje por imitación y el aprendizaje por refuerzo}

En la práctica, es común usar BC como warm-start y luego aplicar RL fine-tune, formando un hybrid pipeline: primero se aprende una policy inicial con demonstrations de alta calidad, y después se mejora la robustez mediante rollout distribuido + reward fine-tuning.

\begin{warningbox}{No abandonar la señal de reward}
Aunque el objetivo sea imitar al experto, también conviene monitorear de manera continua el reward potencial (por ejemplo, safety penalty, human critique score), para poder hacer un rollback de inmediato en cuanto la recompensa presente drift.
\end{warningbox}

\subsection{Resumen de la sección}

Los compounding errors requieren combinar la recolección de datos con estrategias de human-in-the-loop. DAgger/HG-DAgger ofrece una ruta de corrección, y al combinarla con un RL fine-tuning posterior, la policy puede seguir convergiendo bajo distribution shift.

